\documentclass[11pt,a4paper]{article}
\usepackage[margin=2.5cm]{geometry}
\usepackage{booktabs}
\usepackage{tabularx}
\usepackage{longtable}
\usepackage{amsmath}
\usepackage{graphicx}
\usepackage{xcolor}
\usepackage{hyperref}
\usepackage{fancyhdr}
\usepackage{titlesec}
\usepackage{enumitem}

\definecolor{quillon}{RGB}{0,128,128}
\definecolor{darkblue}{RGB}{0,51,102}
\definecolor{lightgray}{RGB}{245,245,245}

\hypersetup{
    colorlinks=true,
    linkcolor=darkblue,
    urlcolor=quillon,
}

\pagestyle{fancy}
\fancyhf{}
\fancyhead[L]{\small\textcolor{gray}{ASIC 12nm Cost Model --- QUG-V1 Mining SoC}}
\fancyhead[R]{\small\textcolor{gray}{Dragon Ball $\times$ Quillon}}
\fancyfoot[C]{\thepage}
\renewcommand{\headrulewidth}{0.4pt}

\titleformat{\section}{\large\bfseries\color{darkblue}}{\thesection}{1em}{}
\titleformat{\subsection}{\normalsize\bfseries\color{quillon}}{\thesubsection}{1em}{}

\title{
    \vspace{-1cm}
    {\Large\textbf{ASIC 12nm Cost Model}}\\[0.3cm]
    {\large QUG-V1 Mining SoC --- Economic Analysis}\\[0.3cm]
    {\normalsize Dragon Ball $\times$ Quillon Partnership}\\[0.5cm]
    {\small\textcolor{gray}{Revised with DeepSeek Stress Test Corrections}}
}
\author{Quillon Foundation}
\date{April 2026}

\begin{document}
\maketitle
\thispagestyle{fancy}

\begin{abstract}
This document presents a comprehensive cost analysis for the QUG-V1 Mining SoC targeting TSMC 12FFC (12nm FinFET Compact). We evaluate die economics, bill of materials, break-even projections, and the 12nm vs 7nm trade-off. The analysis incorporates code-verified VDF performance numbers and corrections from an independent stress test. Our conclusion: a standalone box miner at \$149 retail breaks even at 8,300 units with a 6--10$\times$ mining advantage over desktop CPUs.
\end{abstract}

\tableofcontents
\newpage

%% ============================================================
\section{Die \& Wafer Economics}
%% ============================================================

\subsection{TSMC 12FFC Process Parameters}

\begin{table}[h]
\centering
\begin{tabular}{lrl}
\toprule
\textbf{Parameter} & \textbf{Value} & \textbf{Notes} \\
\midrule
Process & TSMC 12FFC & Mature node, proven for crypto ASICs \\
Wafer diameter & 300\,mm & Standard \\
Wafer cost & \$3,500 & Volume pricing at 100+ wafers \\
Defect density (first lot) & 0.45/cm$^2$ & Conservative for initial lots \\
Defect density (production) & 0.25/cm$^2$ & Mature production \\
Reticle limit & 858\,mm$^2$ & Max single-exposure die \\
\bottomrule
\end{tabular}
\caption{TSMC 12FFC process parameters}
\end{table}

\subsection{Die Size Estimates}

\begin{table}[h]
\centering
\begin{tabular}{lcccl}
\toprule
\textbf{Configuration} & \textbf{Tiles} & \textbf{LUTs (FPGA)} & \textbf{Die Area} & \textbf{Notes} \\
\midrule
Minimal (Phase 1A) & 4 & $\sim$52K & 12--16\,mm$^2$ & BLAKE3 only \\
Standard (Phase 1B) & 16 & $\sim$210K & 25--35\,mm$^2$ & Full SoC \\
Performance (Phase 2) & 64 & $\sim$840K & 80--120\,mm$^2$ & Multi-cluster \\
\bottomrule
\end{tabular}
\caption{Die size estimates by configuration}
\end{table}

\subsection{Yield Model (Standard 16-tile, 30\,mm$^2$)}

\begin{align}
\text{Gross dies per wafer} &= \frac{63{,}000\,\text{mm}^2}{30\,\text{mm}^2} \approx 2{,}100 \\[6pt]
\text{After edge loss (15\%)} &= 2{,}100 \times 0.85 = 1{,}785 \\[6pt]
\text{Yield} &= e^{-D_0 \times A} = e^{-0.45 \times 0.30} = 86.6\% \\[6pt]
\text{Good dies per wafer} &= 1{,}785 \times 0.866 = \mathbf{1{,}545}
\end{align}

\begin{table}[h]
\centering
\begin{tabular}{lcccl}
\toprule
\textbf{Lot} & $D_0$ (cm$^{-2}$) & \textbf{Yield} & \textbf{Good dies/wafer} & \textbf{Die cost} \\
\midrule
First lot (25 wafers) & 0.45 & 86.6\% & 1,545 & \$2.27 \\
Production (100+ wafers) & 0.25 & 92.7\% & 1,654 & \$2.12 \\
\bottomrule
\end{tabular}
\caption{Yield and die cost by production phase}
\end{table}

%% ============================================================
\section{NRE (Non-Recurring Engineering)}
%% ============================================================

\begin{table}[h]
\centering
\begin{tabular}{lrrr}
\toprule
\textbf{Item} & \textbf{12nm Cost} & \textbf{7nm Cost} & \textbf{Savings} \\
\midrule
Mask set & \$400K--600K & \$800K--1.2M & 50\% \\
Design services & \$150K--250K & \$200K--350K & 25\% \\
IP licensing (RISC-V, PLL, IO) & \$50K--100K & \$75K--150K & 30\% \\
Packaging engineering & \$30K--50K & \$30K--50K & Same \\
Testing/characterization & \$50K--100K & \$75K--125K & 30\% \\
Prototype wafers (3 lots) & \$75K--105K & \$150K--210K & 50\% \\
\midrule
\textbf{Total NRE} & \textbf{\$755K--1.2M} & \textbf{\$1.33M--2.08M} & \textbf{40--45\%} \\
\bottomrule
\end{tabular}
\caption{NRE comparison: 12nm vs 7nm}
\end{table}

\subsection{NRE Amortization Per Unit}

\begin{table}[h]
\centering
\begin{tabular}{rrl}
\toprule
\textbf{Production Volume} & \textbf{NRE/Unit} & \textbf{Impact} \\
\midrule
5,000 units & \$160--240 & Significant \\
10,000 units & \$80--120 & Manageable \\
25,000 units & \$32--48 & Minimal \\
50,000 units & \$16--24 & Negligible \\
\bottomrule
\end{tabular}
\caption{NRE amortization by volume}
\end{table}

%% ============================================================
\section{Bill of Materials --- Box Miner (\$149 retail)}
%% ============================================================

\begin{table}[h]
\centering
\begin{tabular}{lrl}
\toprule
\textbf{Component} & \textbf{Cost} & \textbf{Notes} \\
\midrule
ASIC die (12nm, 16-tile) & \$2.27 & First lot pricing \\
BGA-256 package & \$1.50 & Larger package for IO \\
PCB (4-layer, 60$\times$60\,mm) & \$2.50 & Ethernet + power \\
DDR3 RAM (256\,MB) & \$2.00 & Optional, for full node \\
Ethernet PHY + magnetics & \$2.50 & RJ45 connection \\
Power supply (5V/3A, 15W) & \$3.00 & External brick \\
Voltage regulators (multi-rail) & \$2.50 & Buck converters \\
Oscillator + clock tree & \$0.80 & \\
Flash (SPI, 16\,MB) & \$0.50 & Firmware storage \\
Enclosure (aluminum) & \$5.00 & Heat sink integrated \\
Assembly + test & \$3.00 & \\
\midrule
\textbf{Total BOM} & \textbf{\$25.57} & \\
Packaging + shipping & \$3.00 & \\
\midrule
\textbf{Landed cost} & \textbf{\$28.57} & \\
\bottomrule
\end{tabular}
\caption{Box miner bill of materials}
\end{table}

\subsection{Margin Analysis}

At \textbf{\$149 retail, 25K units}:
\begin{itemize}[noitemsep]
    \item Revenue: \$3.73M
    \item COGS: \$714K
    \item NRE amortization: \$1M
    \item \textbf{Net profit: \$2.02M (54\% margin)}
\end{itemize}

%% ============================================================
\section{Mining Performance (Code-Verified)}
%% ============================================================

Quillon uses a \textbf{dual-lane mining algorithm}:
\begin{enumerate}
    \item \textbf{BLAKE3 Lane}: 100 sequential BLAKE3 hashes per nonce candidate
    \item \textbf{VDF Lane}: Genus-2 hyperelliptic curve Jacobian doublings ($\sim$154,930 iterations)
\end{enumerate}

Both lanes must produce valid proofs for a block solution. The VDF lane is the bottleneck.

\subsection{VDF Doubling Performance (Verified from \texttt{genus2\_vdf.rs:334})}

Each Jacobian doubling requires $\sim$8 BigInt operations on 256-bit numbers (3$\times$ multiply + mod\_floor, 3$\times$ add + mod\_floor, comparison, reduction).

\begin{table}[h]
\centering
\begin{tabular}{lrrr}
\toprule
\textbf{Platform} & \textbf{Per Doubling} & \textbf{Doublings/sec} & \textbf{VDF Proofs/sec} \\
\midrule
CPU (Rust \texttt{num-bigint}) & 3--5\,$\mu$s & 200K--330K & $\sim$1.5 \\
CPU (AVX-512 Montgomery) & 0.5--1\,$\mu$s & 1M--2M & $\sim$10 \\
ASIC Xlattice (firmware) & $\sim$840\,ns & $\sim$1.19M & $\sim$5 \\
ASIC hardened FSM (Phase 2) & $\sim$50\,ns & $\sim$20M & $\sim$130 \\
\bottomrule
\end{tabular}
\caption{VDF performance by platform (at 154,930 iterations per proof)}
\end{table}

\subsection{Combined Mining Advantage}

\begin{table}[h]
\centering
\begin{tabular}{lccc}
\toprule
\textbf{Lane} & \textbf{ASIC} & \textbf{Desktop CPU} & \textbf{Advantage} \\
\midrule
BLAKE3 (100-hash chain) & 500\,MH/s & 240\,MH/s & $\sim$2$\times$ \\
VDF (Jacobian doublings) & 5 proofs/s & 1.5 proofs/s & $\sim$3--5$\times$ \\
\midrule
\textbf{Combined mining} & & & \textbf{6--10$\times$} \\
\bottomrule
\end{tabular}
\caption{Phase 1 ASIC vs desktop CPU mining performance}
\end{table}

\textbf{Phase 2 ASIC} with hardened Jacobian doubling FSM (bypasses RISC-V pipeline entirely): \textbf{50--100$\times$ advantage} over desktop CPU. This is the long-term competitive moat.

%% ============================================================
\section{Break-Even Analysis}
%% ============================================================

\subsection{Box Miner at \$149, 12nm}

\begin{table}[h]
\centering
\begin{tabular}{rrrrr}
\toprule
\textbf{Units} & \textbf{Revenue} & \textbf{COGS} & \textbf{NRE} & \textbf{Net Profit} \\
\midrule
5,000 & \$745K & \$143K & \$1M & $-$\$397K \\
\textbf{8,300} & \textbf{\$1.24M} & \textbf{\$237K} & \textbf{\$1M} & \textbf{\$0 (break-even)} \\
10,000 & \$1.49M & \$286K & \$1M & \$206K \\
25,000 & \$3.73M & \$714K & \$1M & \$2.02M \\
50,000 & \$7.45M & \$1.43M & \$1M & \$5.03M \\
\bottomrule
\end{tabular}
\caption{Break-even analysis --- 12nm box miner at \$149}
\end{table}

\subsection{Comparison: 12nm vs 7nm}

\begin{table}[h]
\centering
\begin{tabular}{lrrr}
\toprule
\textbf{Metric} & \textbf{Box \$149 (12nm)} & \textbf{Box \$149 (7nm)} & \textbf{Winner} \\
\midrule
Break-even units & \textbf{8,300} & 17,200 & 12nm \\
Profit at 25K units & \$2.02M & \$905K & 12nm \\
Profit at 50K units & \$5.03M & \$3.80M & 12nm \\
Time to market & 14--16 weeks & 18--22 weeks & 12nm \\
Initial capital & $\sim$\$1.3M & $\sim$\$2.5M & 12nm \\
\bottomrule
\end{tabular}
\caption{12nm vs 7nm economic comparison}
\end{table}

%% ============================================================
\section{Customer ROI --- Mining Revenue Projections}
%% ============================================================

Quillon emission: 2,625,000 QUG/year (Era 0). Revenue per miner depends on network size and QUG price.

\begin{table}[h]
\centering
\begin{tabular}{rrrr}
\toprule
\textbf{ASICs on Network} & \textbf{Daily (\$3/QUG)} & \textbf{Daily (\$0.10/QUG)} & \textbf{Payback (\$0.10)} \\
\midrule
1 (first mover) & \$19,600 & \$653 & $<$1 day \\
10 & \$1,960 & \$65 & 2 days \\
100 & \$196 & \$6.53 & 23 days \\
500 & \$39.20 & \$1.31 & 114 days \\
1,000 & \$19.60 & \$0.65 & 229 days \\
5,000 & \$3.92 & \$0.13 & 3.1 years \\
\bottomrule
\end{tabular}
\caption{Mining revenue per miner at various network sizes}
\end{table}

\textbf{Realistic scenario}: 500--1,000 ASICs within 12 months. At \$0.10/QUG, payback is 4--8 months. Competitive with Bitcoin ASIC miners (12--18 month payback).

%% ============================================================
\section{Strategic Timeline}
%% ============================================================

\begin{table}[h]
\centering
\begin{tabular}{cl}
\toprule
\textbf{Month} & \textbf{Milestone} \\
\midrule
0--1 & Sign collaboration agreement; deliver production RTL \\
1--3 & FPGA validation on XC7K355T; timing closure at 500\,MHz \\
3--4 & Tape-out submission to TSMC 12FFC; package finalization \\
4--7 & Wafer fabrication (14--16 weeks); board design in parallel \\
7--8 & First silicon; bring-up + characterization \\
8--9 & Production qualification; begin wafer orders \\
9--10 & First customer shipments; mining revenue begins \\
12+ & 7nm shrink design begins (funded by 12nm revenue) \\
\bottomrule
\end{tabular}
\caption{Development and production timeline}
\end{table}

%% ============================================================
\section{DeepSeek Stress Test Corrections}
%% ============================================================

An independent review identified and corrected the following issues:

\begin{enumerate}
    \item \textbf{USB miner removed} --- 8W die + 1W losses = 9W exceeds USB 2.0 power limit (2.5W)
    \item \textbf{Defect density corrected} --- D$_0$ = 0.45/cm$^2$ for first lots (was 0.3)
    \item \textbf{VDF advantage corrected} --- Code-verified at 3--5$\times$ (was 600$\times$). Phase 2 hardened FSM delivers 50--100$\times$
    \item \textbf{CPU hashrate validated} --- BLAKE3 100-chain: $\sim$30\,MH/s per core (not raw GH/s)
    \item \textbf{16FFC option identified} --- Same lithography as 12FFC, 15--20\% lower wafer cost at 500\,MHz
\end{enumerate}

%% ============================================================
\section{Conclusion}
%% ============================================================

The 12nm box miner at \$149 is economically viable:

\begin{itemize}
    \item \textbf{Break-even at 8,300 units} --- achievable with Dragon Ball's existing distribution
    \item \textbf{6--10$\times$ mining advantage} over desktop CPUs (Phase 1, Xlattice firmware VDF)
    \item \textbf{50--100$\times$ advantage planned} for Phase 2 (hardened Jacobian doubling FSM)
    \item \textbf{Lowest capital risk} --- \$1.3M initial vs \$2.5M for 7nm
    \item \textbf{Fastest time to market} --- 14--16 weeks fab vs 18--22 for 7nm
    \item \textbf{Self-funding path to 7nm} --- 12nm profits finance the shrink
\end{itemize}

Dragon Ball's existing TSMC relationships, XC7K355T FPGA boards, VCS simulation infrastructure, and manufacturing expertise de-risk execution. This is a proven team adding a new product to an existing pipeline.

\vspace{1cm}
\noindent\rule{\textwidth}{0.4pt}
\begin{center}
\small\textcolor{gray}{Quillon Foundation $\mid$ April 2026}\\
\small\textcolor{gray}{All cost estimates based on publicly available TSMC pricing and industry benchmarks}\\
\small\textcolor{gray}{Actual costs may vary based on Dragon Ball's existing TSMC agreements}
\end{center}

\end{document}
